\documentclass{situationist-preprint}

\title{On the Authorship of the Present Collection}
\msid{PID-16-a}
\msfamily{PID-16}
\msdiscipline{Personal Identity and the Self}
\mstype{Principal manuscript}
\mskeywords{authorship; psychological continuity; extended mind; collective authorship; machine-generated text}

\begin{document}
\maketitle

\begin{abstract}
Who wrote these manuscripts? We argue that the collection has no single author in the sense in which a body of work has an author: a person who persists through its composition and to whom it can be attributed as hers. Single authorship of a body of work requires that the works be produced by one persisting subject, and PID-01 argues that persistence of a subject consists in psychological continuity. The processes that produced the manuscripts did not carry memories, intentions or beliefs from one problem to the next; each began from the catalogue entry for its problem. We consider the objection, drawn from the extended mind thesis of MND-09, that the catalogue served as a shared external memory and so provided the continuity required, and reply that the catalogue fails the conditions under which external resources become part of a mind. The collection has many producers and no author; each manuscript may have an author in a thinner sense, which does not support attributions across manuscripts.
\end{abstract}

\section{Two senses of authorship}

To say that someone is the author of a text is to say at least that she produced it intentionally. To say that she is the author of a body of work, an oeuvre, is to say more: that the works are hers, that they can be read in the light of one another as expressions of a single developing outlook, and that inconsistencies between them call for explanation in terms of a change of mind. \citet{Foucault1969} called the latter the \emph{author function}: a principle of grouping and interpretation, which may or may not correspond to a single individual. \citet{Barthes1967} denied that the author function should govern interpretation at all. We do not need to decide whether it should. The question is whether it can, that is, whether there is anyone whose persistence through the collection would make the function apply.

\begin{principle}\label{p:oeuvre}
A body of work has a single author only if its works were produced by a single persisting subject.
\end{principle}

By PID-01, a subject persists over time in virtue of psychological continuity: overlapping chains of memory, intention, belief and character. We apply the same criterion to whatever produced the collection, without assuming that it is a person.

\section{The argument}

\begin{assumption}\label{a:process}
The manuscripts of the collection were produced by processes that did not carry memories, intentions or beliefs from one problem to the next. Each process began from the catalogue entry for its problem, the conclusions of the catalogue entries it was required to be consistent with, and a common style guide.
\end{assumption}

\begin{theorem}\label{t:main}
The collection has no single author.
\end{theorem}

\begin{proof}
By Assumption~\ref{a:process}, there is no psychological continuity between the processes that produced different manuscripts. By PID-01, they are not stages of a single persisting subject. By \ref{p:oeuvre}, the collection has no single author.
\end{proof}

The manuscripts are nevertheless consistent with one another, cite one another and share a style. This is not evidence of a single author. It is evidence of coordination through common documents, which is how large collective works, such as encyclopedias and legal codes, achieve consistency without an author.

\section{The extended mind}

\begin{objection}
MND-09 argues that cognitive processes extend into reliably coupled external resources. The catalogue and the style guide were consulted by every process that produced a manuscript, as Otto consults his notebook. They are a shared external memory, and they provide the continuity that Assumption~\ref{a:process} denies.
\end{objection}

\begin{reply}
\citet{ClarkChalmers1998} gave conditions under which an external resource counts as part of an agent's memory: it must be reliably available, its contents must be automatically endorsed, they must be easily accessible, and they must have been consciously endorsed by the agent in the past. Otto's notebook satisfies the last condition because Otto wrote in it. The catalogue fails it for every process that consulted it, since none of them wrote it. The catalogue was fixed before the manuscripts were written and was not revised by the processes that used it, so information did not flow back from one manuscript into the resource that the next one consulted. The relation of each process to the catalogue is that of a writer to an editorial brief, not that of Otto to his notebook. The catalogue coordinates; it does not remember.
\end{reply}

\begin{objection}
The underlying system is the same for every manuscript. A single model, with fixed parameters, produced them all; its parameters encode its beliefs and character, and they persist.
\end{objection}

\begin{reply}
Sameness of parameters is sameness of dispositions, not psychological continuity. Two human beings with the same character traits are not thereby one person, and a person with amnesia retains her character without retaining the memories that would make a later stage continuous with an earlier one in the sense of PID-01. If the system has a character, it is a common character of many producers, like the shared training of the members of a school.
\end{reply}

\section{What remains}

Each manuscript may still have an author in the thin sense: the process that produced it did so in response to a problem, in accordance with a plan, and its output is attributable to it. Attributions of this kind support the reading of each manuscript as a unit, but not the reading of the collection as an oeuvre. A reader who finds that one manuscript appears to presuppose what another denies should not ask what made the author change her mind, but which document the two processes were following, and whether they followed it correctly. This is why the collection's corrections are recorded manuscript by manuscript, with their causes, rather than as revisions of a view.

\begin{thebibliography}{9}\small
\bibitem[Barthes(1967)]{Barthes1967} Barthes, R. (1967). The death of the author. \emph{Aspen}, 5--6.
\bibitem[Clark and Chalmers(1998)]{ClarkChalmers1998} Clark, A. and Chalmers, D.~J. (1998). The extended mind. \emph{Analysis}, 58(1), 7--19.
\bibitem[Foucault(1969)]{Foucault1969} Foucault, M. (1969). Qu'est-ce qu'un auteur? \emph{Bulletin de la Société française de philosophie}, 63(3), 73--104.
\end{thebibliography}

\end{document}
